\documentclass{article}
\usepackage[utf8]{inputenc}
\usepackage{amsmath}
\usepackage{geometry}
\geometry{margin=2cm}
\begin{document}

\section*{Resolução da Questão 131 — ENEM 2024 — Ciências da Natureza}

\subsection*{Tema: Sucessão ecológica e mudanças em comunidades vegetais}

\paragraph{Interpretação do enunciado}
A questão apresenta um texto descrevendo sistemas agroflorestais que envolvem a introdução progressiva de espécies vegetais, causando mudanças graduais na composição e estrutura das comunidades vegetais, culminando no aumento da diversidade do ambiente. O desafio é identificar qual processo natural é análogo a essas mudanças graduais.

\paragraph{Termos-chave}
\begin{itemize}
  \item \textbf{Análogo:} comparar um conjunto de mudanças graduais a um processo natural conhecido.
  \item \textbf{Processo natural:} fenômeno que ocorre na natureza sem intervenção direta humana.
\end{itemize}

\paragraph{O que a questão pede}
Identificar, entre as opções, o processo natural que corresponde ao fenômeno descrito no texto: mudanças graduais e aumento da diversidade em comunidades vegetais.

\paragraph{Estratégia para resolução}
\begin{itemize}
  \item Interpretar que o enunciado descreve mudanças graduais e o aumento de diversidade na vegetação.
  \item Comparar essas características com conceitos ecológicos das alternativas.
  \item Descartar termos que não retratam mudanças graduais naturais em comunidades (rot. de culturas, coevolução, adaptação gen., convergência).
  \item Confirmar que a sucessão ecológica define esse processo natural gradativo.
\end{itemize}

\subsection*{Revisão conceitual dos principais termos}

\begin{tabular}{|p{3cm}|p{6cm}|p{6cm}|}
\hline
\textbf{Conceito} & \textbf{Ideia-chave} & \textbf{Aplicação ao texto da questão} \\ \hline
Sucessão Ecológica & Mudanças graduais na composição e estrutura das comunidades ao longo do tempo & Descreve as mudanças graduais que aumentam a diversidade ambiental mencionadas no enunciado \\ \hline
Rotação de Culturas & Alternância planejada de culturas agrícolas para manejo do solo & Prática agrícola, não processo natural; citado no texto, mas não o análogo buscado \\ \hline
Coevolução Específica & Evolução conjunta entre espécies por interação biológica & Não corresponde a mudança progressiva comunitária descrita \\ \hline
Adaptação por Seleção & Mudança genética em populações induzida pela seleção natural & Processo evolutivo genético, não mudança de composição comunitária \\ \hline
Convergência Adaptativa & Evolução independente de características semelhantes em espécies distintas & Não envolve mudança gradual da comunidade \\ \hline
\end{tabular}

\subsection*{Fundamentos teóricos}

A sucessão ecológica é um processo natural em que a comunidade biológica evolui gradativamente, com alterações na estrutura e composição das espécies, frequentemente aumentando a diversidade e complexidade do ecossistema. Exemplos ocorrem em áreas que sofreram perturbações e estão em recuperação.

É fundamental distinguir a sucessão ecológica das outras opções: rotação de culturas é manejo agrícola; coevolução refere-se à evolução conjunta entre duas ou mais espécies; adaptação por seleção é um processo genético dentro das populações; e convergência adaptativa é uma evolução convergente de características, não da comunidade como um todo.

\subsection*{Resolução final}

Ao analisar o texto do enunciado:
\begin{itemize}
  \item Identificamos a progressão gradual na comunidade vegetal, pela adição de espécies conforme um arranjo espacial e temporal.
  \item Reconhecemos o aumento da diversidade como um resultado natural dessas mudanças.
  \item Confirmamos que o fenômeno descrito é a sucessão ecológica, processo natural de transformação gradual das comunidades.
  \item Eliminamos as demais alternativas por não corresponderem ao processo comunitário gradual descrito.
\end{itemize}

Assim, a alternativa correta é \textbf{B — sucessão ecológica}.

\paragraph{Pulo do gato}
Em provas, lembre que \textit{sucessão ecológica} é o termo-chave para processos naturais que envolvem mudanças graduais e aumento da diversidade em comunidades vegetais, especialmente em áreas em recuperação.

\subsection*{Armadilhas e falsas alternativas}

\begin{itemize}
  \item \textbf{Alternativa A (Rotação de culturas):} Parece válida por mencionar culturas agrícolas no texto, mas trata-se de prática agrícola, diferente de processo natural.
  \item \textbf{Alternativa C (Coevolução específica):} Relacionada à interação evolutiva entre espécies, não a mudança gradual na comunidade.
  \item \textbf{Alternativa D (Adaptação por seleção):} Processo evolutivo genético, não de transformação comunitária vegetal.
  \item \textbf{Alternativa E (Convergência adaptativa):} Evolução independente de características em espécies diferentes, não mudança da comunidade.
\end{itemize}

\subsection*{Código da Questão}

CN_SUCESSAO_ECOLOGICA_001


\end{document}